\documentclass{nfr}

% This LaTeX template is based on the official Word template found here: https://www.forskningsradet.no/siteassets/utlysninger/vedlegg-utlysninger/template-for-project-description-researcher-project-for-experienced-scientists-and-early-career-scientists.docx
% See the README.md file for more information

% Here's an example bibliography using the numeric (Nature) style and two
% columns for the reference list.

% \usepackage{multicol}
% \renewcommand{\bibpreamble}{\begin{multicols*}{2}}
% \renewcommand{\bibpostamble}{\end{multicols*}}

% \usepackage{biblatex}
% \addbibresource{references1.bib}

% \citestyle{nature}
% \citestyle{}
\usepackage{comment}

% \bibliographystyle{bst/veryshortnum}
\bibliographystyle{apalike}

% \title{Research Statement}
\title{Application for Assistant Professor at Reykjavik University}
\author{Zhongli Wang}

\begin{document}

% Project description template
% Researcher Project for Experienced Scientists/Early Career Scientists

% This template consists of
% 
% \textbf{A:} Overview of the chapters and sections to be used as the
% structure of the project description
% 
% \textbf{B:} General guidance
% 
% \textbf{C:} Guidance on the content in each chapter and section
% 
% \textbf{A: Chapters and sections to be used}

% I used pandoc to translate the word template to tex. The formatting of the
% appendices are, therefore, not great, but the contents are there.

\maketitle

% Understanding firm competition in modern markets remains a fundamental challenge in economics. 
% While classical models such as Bertrand, Cournot, and Hotelling provide valuable benchmarks, real markets are characterized by firm asymmetries, network effects, and heterogeneous consumers, which complicate predictions and strategic decision-making \citep{Hotelling1929}.
% My research tackles this challenge by combining computational and statistical modeling to study strategic behavior and market dynamics in complex economic environments. Using tools from agent-based modeling, industrial organization, and econometrics, I analyze how firms and consumers interact in digital and networked markets. The broader goal is to bridge theory, simulation, and data to generate solutions that improve market understanding and support effective policy and managerial decisions.

% \section*{Research Statement}
\noindent Dear Members of the Recruitment Committee,\\[0.2cm]

I am writing to apply for the Assistant Professor position in Management in the Department of Business and Economics at Reykjavik University, commencing in 2026.
I am currently completing my Ph.D. in Economics at Bielefeld University and Università Cattolica del Sacro Cuore within the Marie Skłodowska-Curie EPOC program. 
My research examines how firms make strategic decisions in technology-intensive and rapidly evolving environments, with a particular focus on digital transformation, innovation, and the organizational consequences of regulation and market design. 
Using computational modeling, simulation-based methods, and empirical analysis, I study how firms adapt their capabilities, structure incentives, and respond to uncertainty in dynamic and competitive contexts. 
These themes connect directly to contemporary debates in strategic management and organizational research, and align strongly with Reykjavik University’s commitment to producing rigorous, relevant, and internationally oriented scholarship.\\

I began my PhD in 2021, at a moment when digitalization, remote services, and data governance were reshaping how firms compete and organize their activities. 
This environment motivated the central chapter of my dissertation, which examines how firms in technology-intensive markets make strategic capacity and pricing decisions under network effects, digital infrastructures, and evolving institutional requirements such as data-storage regulations. 
To analyze these dynamics, I combine firm-level data with simulation-based and computational methods to assess how alternative regulatory or organizational choices influence competitive outcomes, strategic incentives, and long-term investment behavior. 
Across its three single-authored chapters, my dissertation integrates theory, quantitative modeling, and computational analysis to show how digital architectures, information frictions, and institutional constraints shape market structure, strategic behavior, and organizational adaptation in rapidly changing industries.\\

Beyond my dissertation, my coauthored papers on \textit{Polytope Fraud Theory} and \textit{Agentic Workflows for Economic Research} extend my interests toward corporate behavior, organizational integrity, and the strategic implications of artificial intelligence \citep{Zhao2025, Dawid2025}. 
The first paper applies machine-learning and statistical techniques to detect patterns of corporate fraud, offering insights into how firms manage transparency, governance, and stakeholder trust. 
The second develops a structured framework for integrating large-language-model agents into empirical research workflows, illustrating how AI can reshape analytical capabilities, knowledge production, and organizational decision-making. 
Together, these projects reflect my commitment to combining rigorous analytical methods with questions of strategic relevance—such as responsible data practices, technological innovation, and the evolving relationship between organizations, digital tools, and society—topics that align well with Reykjavik University's emphasis on impactful and interdisciplinary management research.\\

At Reykjavik University, I hope to contribute my theoretical, empirical, and computational expertise to research on strategic decision-making in digital and technology-intensive environments. 
I am particularly interested in how firms navigate information frictions, regulatory pressures, and technological change, and how these forces shape competitive dynamics, organizational capabilities, and long-term value creation. 
Building on my experience with simulation-based modeling, firm-level data, and computational methods, I aim to develop projects that examine strategic adaptation in digital ecosystems, the diffusion of technological shocks across networks of firms and stakeholders, and the organizational and societal implications of emerging data and AI practices.\\
% Such work would complement the Department’s strengths in strategic management, corporate sustainability, and business \& society, and aligns with IÉSEG’s mission to promote responsible, impactful research that prepares organizations and decision-makers for a rapidly changing world.\\
% More broadly, my research interests include:
% \begin{itemize}
%     \item Strategic behavior and organizational adaptation in digital, data-intensive, and technology-driven environments;
%     \item Computational and theoretical approaches to understanding how firms develop capabilities, make decisions under uncertainty, and sustain competitive advantage;
%     \item Simulation-based and machine-learning methods for analyzing strategic interactions, technological change, and the diffusion of innovations across firm and stakeholder networks;
%     \item The role of digital infrastructures, algorithms, and data governance in shaping organizational practices, platform ecosystems, and broader societal outcomes.
% \end{itemize}
% I am particularly motivated by the opportunity to contribute to IÉSEG’s Management \& Society Department, whose commitment to interdisciplinary, responsible, and impact-oriented research aligns strongly with my work on strategic decision-making, digital transformation, and the societal implications of emerging technologies.\\

For empirical analysis, I draw on a range of digital, organizational, and market-level datasets, including access to the Refinitiv platform through my collaboration with the Southern University of Science and Technology (SUSTech). 
These data resources allow me to examine how firms operate in technology-intensive environments, how strategic decisions evolve under uncertainty, and how digital infrastructures and regulatory contexts shape organizational behavior and stakeholder outcomes. 
My doctoral research, funded by the \textit{Marie Skłodowska-Curie Actions} and the \textit{Bielefeld Young Researchers’ Fund}, has demonstrated my ability to manage externally funded projects, collaborate across international institutions, and coordinate large-scale empirical and computational research initiatives. 
This combination of analytical rigor, international collaboration, and cross-disciplinary engagement positions me to contribute effectively to Reykjavik University’s research community and its broader commitment to producing impactful and practice-relevant management scholarship.\\

I am deeply committed to teaching, and my philosophy centers on clarity, inclusion, and active engagement. 
I view teaching as a process of joint exploration, and I aim to build students’ analytical confidence by linking concepts to real-world managerial challenges through data-driven examples, case discussions, and hands-on computational exercises. 
During my master’s studies at the University of Zurich, I served for two years as a teaching assistant and contributed to the SYLVA project, an automated scoring and feedback system that enhanced learning effectiveness and fairness in large courses. 
This experience strengthened my commitment to integrating technology with human-centered pedagogy. 
More recently, I co-delivered a public lecture in German on “AI Agents in Science,” which reinforced the importance of clarity and adaptability when communicating complex ideas. 
I am prepared to teach courses in strategic management, digital transformation, and analytical methods at the undergraduate and graduate levels, and I would welcome the opportunity to contribute to Reykjavik University’s student-centered learning environment by helping students develop the analytical and critical-thinking skills needed to navigate technology-driven organizational contexts.\\

Going forward, I plan to develop a focused research agenda that integrates computational and empirical methods to investigate strategic decision-making, digital transformation, and the organizational implications of emerging technologies. 
I aim to pursue externally funded interdisciplinary projects that examine how firms adapt to technological change, how algorithmic and data-driven tools reshape strategic capabilities, and how organizations can implement responsible and sustainable practices in digital environments. 
I am particularly interested in collaborating with scholars in strategic management, information systems, and organizational studies, as well as with industry and societal stakeholders, to translate analytical insights into actionable frameworks that enhance strategic resilience, ethical data governance, and innovation.\\
% With IÉSEG’s strong institutional support, international research culture, and commitment to responsible management, I am confident that this research agenda will contribute meaningfully to the School’s mission of empowering changemakers and producing impactful scholarship for academic and managerial communities.\\

Thank you for considering my application! 
I would be honored to contribute to Reykjavik University’s mission by supporting high-quality research and delivering innovative, engaging teaching.\\
\vspace{5mm}

\noindent Sincerely,\\
\vspace{5mm}

\noindent Zhongli Wang

\vspace{5mm}

% I would be delighted to join IÉSEG School of Management and contribute to its vibrant, international, and values-driven academic community. 
% I am deeply motivated by the School’s commitment to responsible management, interdisciplinary research, and societal impact, and 


% Throughout my doctoral and collaborative research, I have experienced the value of combining analytical modeling with computational innovation. Working across fields has shown me that the most pressing problems—such as market transparency, sustainable technology adoption, and responsible data governance—can only be solved through interdisciplinary collaboration.

% In addition to my market-oriented research, my coauthored publication \textit{``Polytope Fraud Theory''} and the working paper \textit{``Agentic Workflows for Economic Research''} have shaped my commitment to methodological innovation. The former integrates economic reasoning with machine learning to enhance fraud detection and investor protection; the latter connects computational economics with artificial intelligence by designing LLM-based research workflows. These experiences have strengthened my belief that economists can contribute not only to theoretical insights but also to the design of intelligent, transparent, and reproducible systems.

% I am therefore particularly excited about joining the FINATRAX Research Group, where business research meets information systems engineering. The group’s mission—to explore digital identities, distributed ledgers, machine learning, and IoT as enablers of sustainable technologies—closely aligns with my research vision. I am enthusiastic about participating in the design, implementation, and theoretical elaboration of high-quality, interdisciplinary solutions that create tangible value for both organizations and society. This position represents a natural continuation of my efforts to bridge economics with emerging digital technologies.

% \section*{Funding and Collaboration Strategy}

% My strategy for research funding and partnerships builds on a dual foundation of analytical rigor and applied collaboration. 
% Methodologically, I plan to combine analytical and computational approaches: analytical derivations will establish benchmark propositions, while numerical simulations and agent-based modeling will explore market dynamics and firm behavior under varying degrees of information transparency. This structure naturally supports both theoretical innovation and data-driven applications—key elements for competitive research proposals.

% For empirical validation, I will leverage access to the Refinitiv database through collaboration with the Southern University of Science and Technology (SUSTech), complemented by additional data resources available at the host institution. My doctoral research, funded by the \textit{Marie Skłodowska-Curie Actions} and the \textit{Bielefeld Young Researchers’ Fund}, has already demonstrated my ability to manage funded projects and coordinate cross-institutional collaboration. 

% Going forward, I intend to apply for European and national-level grants such as \textit{Horizon Europe}, \textit{DFG}, and \textit{VolkswagenStiftung}, focusing on projects that integrate computational economics with digital transformation research. On the industry side, I aim to build partnerships with firms in the ICT and FinTech sectors to co-develop decision-support tools and simulation frameworks for market design and policy evaluation. These collaborations would ensure practical relevance, foster technology transfer, and open opportunities for joint funding applications. With adequate computational resources and institutional support, I am confident that these efforts can yield both high-impact publications and meaningful contributions to industry and policy communities.





\begin{comment}

\section*{Statement of Research Interests}

% Here's a citation to \cite{texbook}. Here are citations to \cite{texbook,knuth:1984,lesk:1977}, and to \cite{lesk:1977,texbook,latex2e,latex:companion}.

My doctoral dissertation, ``\textit{Strategic Behavior and Economic Dynamics in ICT Markets: An Agent-based Framework}", investigates how firms in the information and communications technology (ICT) industry adapt under constraints such as data regulations, capacity limits, and consumer network effects. The work builds a bridge between state-space and agent-based models, showing how techniques from dynamic latent-variable systems can be applied to simulate and estimate complex economic environments.\footnote{All three chapters are single-authored papers.}\\

My job market paper, \textit{``Strategic Firm Capacity Selection under Network Effects: A Computational Model of Capacity Investment Dynamics in the ICT Industry''}, investigates how digital firms determine optimal service capacity under the joint influence of consumer congestion and network effects---an interaction that has received limited attention in the existing literature. While traditional industrial organization models emphasize innovation or price competition, many digital firms face a different constraint: their profitability depends critically on infrastructure investment and the ability to maintain service quality. To address this challenge, I combine game-theoretic reasoning with agent-based simulations. The game-theoretic framework captures the multi-stage nature of firms’ decisions on pricing and capacity investment, while the agent-based model represents the market dynamics that emerge when heterogeneous consumers interact through a small-world network and share information within user communities. The results show that strong network effects tend to amplify competitive advantages by concentrating demand, whereas congestion generates counteracting pressures that reduce pricing flexibility and erode profits. Together, these forces explain why digital firms must balance capacity expansion and user experience to sustain long-term profitability. This study is among the first to systematically analyze the joint role of network and congestion effects in the ICT industry, offering new insights into pricing, investment, and regulatory strategies for digital service markets.\\

Building on the similar market context, my previous paper, \textit{``Heterogeneous Consumers, Firm Competition, and the Data Value Chain: An Agent-Based Approach''}, investigates how differences in data-governance intensity influence competition in software markets. The challenge arises because digital firms depend on user data to enhance product quality and pricing decisions, while local governments impose varying levels of data regulation to protect privacy and maintain market fairness. 
To analyze this tension, I develop a hybrid game-theoretic and agent-based model in which governments determine the stringency of data governance, and firms subsequently compete on investment and pricing decisions. 
Consumers are heterogeneous in their sensitivity to product quality and their attitudes toward data governance, which shape their purchasing behavior. 
The results show that stricter data governance limits firms’ ability to utilize consumer data effectively and reduces competitive differentiation, which leads to higher prices and lower consumer welfare. In contrast, moderate levels of regulation can maintain competition and promote more balanced market outcomes. The analysis suggests that regional governments would benefit from greater coordination in setting data policies and from incorporating consumer preferences into policy design, thereby supporting both market efficiency and innovation. Overall, these findings offer a structured perspective on how policymakers can balance privacy protection with incentives for innovation in digital markets.\\

Building on the insights from the first two studies, my third paper ``\textit{Parameter Estimation Methods in an Agent-based Computational Model: Insights on ICT Industrial Dynamics and Policies}" focuses on empirical validation. 
Using financial and operational data from two leading U.S. telecommunication firms, AT\&T and T-Mobile, I estimate the model parameters with the Simulated Method of Moments and neural network–based surrogate modeling (both within-sample and out-of-sample check). 
The main challenge lies in linking complex agent-based models to real-world data, as repeated simulations are computationally intensive and traditional estimation methods often fall short. 
To address this, I employed the Paderborn High-Performance Computing (HPC) clusters and the Python package \texttt{ray} to parallelize the simulations and collect 40{,}000 samples for parameter search. 
By comparing these two approaches, the study provides practical evidence on how to calibrate computational models efficiently while maintaining empirical accuracy. This work contributes to the broader methodological discussion on integrating simulation-based models with data-driven analysis to better understand market behavior and policy outcomes.



\section*{Publications and Emerging Projects}

In addition to my market-oriented research, I have one peer-reviewed publication and one current working paper that contribute to the broader use of computational methods in economics.\footnote{Both papers are coauthored with other researchers.} \\

The paper \textit{``Polytope Fraud Theory''} extends the classical Fraud Triangle and Fraud Diamond frameworks by introducing a multidimensional structure that captures ten abnormal financial practice alarms indicative of potential corporate fraud \citep{Zhao2025}. Using data on activist short sellers and supervised machine learning techniques, the study identifies financial patterns associated with fraudulent reporting and demonstrates how algorithmic approaches can support market transparency and investor protection. 
The paper also demonstrates the potential for automating advanced financial statement analysis, providing tools that can enhance both auditing efficiency and investment decision-making. In addition, it introduces the \textit{Unified Investor Protection Framework}, which systematically categorizes investor-protection theories from the macro- to the micro-level. This framework offers practical guidance for identifying firm-level financial risks and strengthens the connection between analytical modeling and real-world investor decision processes.\\

Another working paper, \textit{``Agentic Workflows for Economic Research: Design and Implementation''}, develops a framework that employs large-language-model (LLM) agents to support economists throughout the research process—from data collection and coding to simulation design and interpretation \citep{Dawid2025}. The framework formalizes how autonomous AI agents can collaborate within multi-stage research pipelines while ensuring reproducibility, version control, and transparent documentation. Using examples implemented in Microsoft AutoGen, we demonstrate how such workflows can accelerate tasks like research-question ideation, literature review, and model specification, thereby improving research efficiency without compromising academic rigor. Importantly, the framework positions human researchers at the core of the process: AI agents act as supportive tools that augment, rather than replace, human judgment and creativity. This line of work extends my methodological interests in AI-enhanced modeling, bridging computational economics with artificial-intelligence research and tool development.

\section*{Future Research}

\subsection*{Potential Research Direction: The Observability of Firm Heterogeneity in Price-Directed Consumer Search}

% Recent models of price-directed search \citep{obradovits2023price} highlight how consumers allocate their attention and search effort across differentiated firms when facing information frictions. However, these frameworks often assume that firm heterogeneity---in cost structures, quality, or service intensity---is either fully observable to consumers or symmetrically distributed across sellers. In many real markets, such as digital platforms, online retailing, or app ecosystems, this heterogeneity is only partially observable: consumers infer firm characteristics indirectly from prices, reputation indicators, or platform signals.

% Understanding how the observability of firm heterogeneity shapes consumer search behavior and market efficiency is crucial for both theory and policy. It determines the extent of equilibrium price dispersion, the efficiency of market matching, and firms’ incentives to disclose information. The problem is empirically feasible because modern datasets (e.g., online pricing panels, platform APIs, or laboratory experiments) provide micro-level evidence on search behavior under varying degrees of information transparency. Computational approaches, including agent-based simulations, further enable the exploration of equilibrium dynamics in complex environments.

Recent models of price-directed search \citep{obradovits2023price} examine how consumers allocate attention and search effort across differentiated firms under information frictions, yet they typically assume that firm heterogeneity---in costs, quality, or service intensity---is fully observable or symmetrically distributed. In real markets such as digital platforms or online retailing, this heterogeneity is only partially observable, as consumers infer firm characteristics indirectly from prices, ratings, or platform signals. Understanding how such observability affects consumer search and market efficiency is crucial, as it shapes equilibrium price dispersion, allocative outcomes, and firms’ incentives to disclose information. The question is empirically feasible given modern micro-level data on online search and pricing, while computational approaches, including agent-based simulations, allow systematic analysis of equilibrium dynamics under varying information transparency. \\

% This research direction contributes to Industrial Organization by linking information structure, firm heterogeneity, and price competition under bounded rationality. It combines theoretical modeling with simulation and empirical validation, making it both conceptually rich and methodologically tractable. 

The following hypotheses will be tested:

\begin{enumerate}%[label=\textbf{H\arabic*.}, leftmargin=2em]
    \item \textbf{Partial observability increases price dispersion.} When firm heterogeneity (e.g., in cost or quality) is imperfectly observable, consumers engage in less directed search, resulting in a wider equilibrium price distribution than under full observability.
    
    \item \textbf{High observability improves allocative efficiency.} Markets with more transparent signals of firm heterogeneity (such as ratings or quality certifications) achieve a closer match between consumers and preferred firms, raising consumer surplus and reducing welfare loss.
    
    \item \textbf{Strategic transparency depends on firm type.} Low-cost or high-quality firms have stronger incentives to increase the observability of their characteristics (e.g., through signaling or advertising) than high-cost or low-quality competitors, leading to asymmetric information equilibria.
\end{enumerate}

\subsection*{Methodological Plan}

In my approach, I plan to combine analytical and computational methods. Analytical derivations will be used to establish the basic propositions and benchmark outcomes of the model, while numerical techniques and computational modeling will be employed to solve for Nash equilibria and explore market dynamics under varying degrees of information transparency. For empirical validation, I have access to the Refinitiv database through an application via the Southern University of Science and Technology (SUSTech) and may also rely on additional database resources at my future host institute. My doctoral research was supported by the Marie Skłodowska-Curie Actions and the Bielefeld Nachwuchsfonds, and it benefited from extensive computational work conducted on Google Colab (self-funded) and the Paderborn High Performance Computing cluster (supported by Bielefeld University). I hope that the new institute could provide adequate support for research infrastructure, particularly computational resources, to further advance this line of study.



% \subsection*{Computational \& ABM Dynamics}

% Specifically, I plan to:

% •	Develop scalable frameworks for data-driven research of economic models that can inform competition and innovation policy.

% •	Study how AI and automation transform firm strategy, market structure, and industrial dynamics.

% My long-term objective is to contribute to the economics literature, that is both analytically rigorous and computationally adaptive—one that captures the rich, evolving behavior of firms and consumers in real markets. I remain deeply motivated to pursue research that connects theoretical innovation with practical relevance, and I look forward to collaborating in an environment that values interdisciplinary approaches and empirical credibility.
% \citep{arve2025decentralized}.


% \bibliography{NFR_Template_Statements/references}

\newpage

\appendix
\end{comment}
\bibliography{NFR_Template_Statements/references1}

\end{document}
